\scanpage{155}
\margin{154}
\uline{Réédition d'E. V. T.\ ,chap.\ II} \mspage{(14)}
\margin{Tribu n\textsuperscript{o} 63}

Remplacer les l.\ 6-11 du bas (d'ailleurs canulées si on n'a pas rectifié l.\ 11) par une récurrence sur $n$: $x=x_1$ ou $x$ barycentre de $x_2,\ldots,x_n$.

P.\ 115, l.\ 4, une référence pour $B(\mathbb{N})$.

Remplacer l'exemple 3 par: soit $u:C\longrightarrow C'$ où $C$ et $C'$ sont convexes et $u$ affine; soient $x$ un point extrémal de $C'$, $y$ un point extrémal de $u^{-1}(x)$; alors $y$ est extrémal dans $C$.

l.\ 6 du bas: compact non vide.

P.\ 116, l.\ 6, (i), (ii) et (iii) sont également triviaux.

P.\ 117, l.\ 1: Puisque $E$ est elc séparé, il suffit de voir \ldots

Vider les cor.\ 1, 2, 3.

P.\ 118, démonstration du th.\ 1: D'après la prop.\ 1, toute fonction affine continue $\geq0$ sur $B$ est $\geq0$ sur $A$. Donc $B=A$ (s'assurer qu'on a bien mis au \S\ de la séparation des convexes exactement ce qu'il faut).

P.\ 118 et 119, échanger la prop.\ 2 et son corollaire. Démonstration du corollaire: Pour tout demi-espace ouvert $D$ contenant $x$, soit $L_D=(E-D)\cap A$ qui est convexe compact. $A-\{x\}$ est convexe, donc $\Gamma(L_D\cup L_{D'})$ ne contient pas $x$, et est compact. Donc les $L_D$ sont filtrant décroissants. Les $L\cap\bar D$ ont pour intersection $x$, et forment un système fondamental de voisinages de $x$. Démonstration de la prop.\ 2: Si $K$ laisse échapper un point extrémal $x$, il laisse échapper une tranche voisinage de $x$, donc $\overline{\Gamma(K)}\neq A$.

Exercice: dans le cor., remplacer $x$ par un chapeau.

P.\ 120, la conjecture de la l.\ 9 est fausse. Voici une contre-exemple (exercice: deviner l'auteur): soit $B$ la boule unité d'un Hilbert faible $H$; $x\in$ sphère unité; il n'y a qu'un hyperplan d'appui; celui-ci disparaît si on affaiblit encore un peu la topologie faible (en utilisant un sous-espace partout dense bien choisi de $H$ ).

l.\ 10 du bas: un cône convexe.

l.\ 8 du bas, ajouter: Si $x\in D$ et $x=y+y'$ ($y,y'\in C$) alors $y,y'\in D$, et ceci caractérise les génératrices extrémales. (La Remarque 1, p.\ 121, devra alors être abrégée).

l.\ 5 du bas: $H$ un hyperplan contenant $x_o$.

P.\ 121, après la déf.\ 3, un exemple: * soient $X$ localement compact, $f$ une fonction
